\documentclass[aspectratio=169,11pt]{beamer}
\usepackage[indonesian]{babel}
\usepackage{tikz,pgfplots,listings}
\usetikzlibrary{positioning,arrows.meta,fit,backgrounds}
\pgfplotsset{compat=1.18}
\usetheme{metropolis}
\setbeamerfont{title}{size=\Large}
\setbeamerfont{author}{size=\small}
\lstset{language={[LaTeX]TeX},basicstyle=\ttfamily\small,columns=fullflexible,keepspaces=true,showstringspaces=false,breaklines=true}
\title{TikZ untuk pengajaran dan penelitian}
\subtitle{Hari 4 --- TikZ Mindmap}
\author{Pelatihan Universitas Terbuka dan Departemen Matematika FMIPA UGM}
\date{}
\begin{document}
\begin{frame}
  \maketitle
\end{frame}
\begin{frame}{Tujuan hari ini}
\begin{enumerate}
\item Sesi 1: membuat gambar dasar dan diagram yang mudah diubah.
\item Sesi 2: grafik akademik, integrasi, dan proyek.
\item Sesi 3: belum ditentukan
\item Sesi 4: peta kuliah dengan TikZ
\end{enumerate}
\end{frame}
\begin{frame}{Empat sesi}
Sesi 1: dasar TikZ dan diagram terstruktur. Sesi 2: grafik akademik, integrasi Beamer, dan proyek akhir.
\medskip

Sesi 3 disiapkan kemudian. Untuk sesi 4, siapkan satu konsep atau tema mata kuliah dari bidang Anda.
\end{frame}
\section{Sesi 1: Dasar dan diagram yang dapat digunakan kembali}
\begin{frame}{Sesi 1: tujuan dan alokasi 75 menit}
\begin{itemize}
\item 0--15: dokumen minimal, koordinat, garis, dan bentuk.
\item 15--25: warna, label, dan latihan dasar singkat.
\item 25--45: node bernama, positioning, gaya, dan scope.
\item 45--60: latihan diagram; fit dan loop sebagai pengayaan.
\item 60--70: diagnosis, refaktor, dan pembahasan.
\item 70--75: exit ticket dan refleksi.
\end{itemize}
Target: membuat ilustrasi dasar dan diagram dengan node bernama serta gaya yang dapat dipakai ulang.
\end{frame}
\begin{frame}{Mengapa menggambar dengan TikZ?}
\begin{itemize}
\item Gambar berbasis vektor: tetap tajam saat diperbesar.
\item Teks dan notasi matematika mengikuti tipografi dokumen.
\item Sumber dapat disunting, dibandingkan, dan digunakan kembali.
\item Tepat untuk diagram konseptual dan ilustrasi yang perlu konsisten.
\end{itemize}
Diskusi singkat: gambar apa yang sering Anda buat untuk kuliah atau penelitian?
\end{frame}
\begin{frame}[fragile]{Dokumen minimal: mulai dari sini}
\begin{lstlisting}
\documentclass{article}
\usepackage{tikz}
\begin{document}
\begin{tikzpicture}
  \draw (0,0) -- (2,1);
\end{tikzpicture}
\end{document}
\end{lstlisting}
Simpan sebagai \texttt{gambar.tex}, lalu kompilasi dengan pdfLaTeX. Periksa PDF dan log jika gagal.
\end{frame}
\begin{frame}[fragile]{Apa yang dilakukan perintah ini?}
\begin{lstlisting}
\draw[blue,thick,->] (0,0) -- (2,1);
\end{lstlisting}
\begin{itemize}
\item \texttt{\textbackslash draw}: gambar path; opsi dalam kurung siku.
\item \texttt{(0,0)} dan \texttt{(2,1)}: titik awal dan akhir.
\item \texttt{--}: segmen garis lurus; titik koma mengakhiri path.
\item \texttt{blue}, \texttt{thick}, dan \texttt{->}: warna, ketebalan, dan panah.
\end{itemize}
\end{frame}
\begin{frame}{Model dasar}
\begin{itemize}
\item Koordinat menentukan posisi; path menghubungkan posisi.
\item Node dapat menempatkan teks atau notasi matematika pada koordinat.
\item Opsi mengatur tampilan; setiap perintah path berakhir dengan titik koma.
\end{itemize}
\begin{tikzpicture}
\draw[->] (0,0) -- (3,0);
\filldraw[fill=blue!15] (1,0) circle (5mm);
\node[above] at (1,0.5) {Konsep};
\end{tikzpicture}
\end{frame}
\begin{frame}[fragile]{Garis, bentuk, dan teks}
\begin{lstlisting}
\usepackage{tikz}  % di preambule
\begin{tikzpicture}
  \draw[->] (0,0) -- (3,0);
  \draw (0,0) circle (5mm);
  \fill[blue!20] (1,0) rectangle (2,1);
  \node[above] at (2,1) {Label};
\end{tikzpicture}
\end{lstlisting}
Ubah satu koordinat, prediksi hasilnya, lalu kompilasi.
\end{frame}
\begin{frame}[fragile]{Sistem koordinat}
\begin{lstlisting}
\draw[help lines,gray!40] (-1,-1) grid (3,2);
\fill (2,1) circle (2pt) node[above] {Kartesius};
\fill (30:2cm) circle (2pt) node[right] {Polar};
\end{lstlisting}
Secara default, satu unit pada sumbu $x$ atau $y$ bernilai 1 cm. Koordinat polar \texttt{(sudut:jarak)} memakai derajat. Grid membantu konstruksi dan dapat dihapus pada hasil akhir.
\end{frame}
\begin{frame}[fragile]{Koordinat dan path}
\begin{lstlisting}
\draw (0,0) -- (2,0) -- (2,1);
\draw (0,0) -- +(1,0) -- +(0,1);
\draw (0,0) -- ++(1,0) -- ++(0,1);
\end{lstlisting}
\texttt{+} relatif tanpa memperbarui posisi saat ini; \texttt{++} memperbaruinya. Perhatikan perbedaan hasil kedua path.
\end{frame}
\begin{frame}[fragile]{Perintah dasar dan apa yang dihasilkan}
\begin{lstlisting}
\draw[blue,thick] (0,0) -- (2,0);
\fill[blue!20] (0,0) circle (4mm);
\filldraw[fill=orange!20,draw=black]
  (1,0) rectangle (2,1);
\node[above] at (1,1) {Teks};
\end{lstlisting}
\texttt{draw} membuat garis tepi; \texttt{fill} mengisi; \texttt{filldraw} melakukan keduanya; \texttt{node} menempatkan isi.
\end{frame}
\begin{frame}[fragile]{Garis, panah, dan opsi}
\begin{lstlisting}
\draw[thin] (0,0) -- (2,0);
\draw[thick,dashed,red] (0,1) -- (2,1);
\draw[<->,blue] (0,2) -- (2,2);
\end{lstlisting}
Opsi memengaruhi path yang mengikutinya. Coba ubah ketebalan, jenis garis, dan arah panah; amati perubahan di PDF.
\end{frame}
\begin{frame}[fragile]{Path: gambar, isi, dan tutup}
\begin{lstlisting}
\draw[thick] (0,0) -- (1,1) -- (2,0);
\fill[blue!20] (0,0) -- (1,1) -- (2,0) -- cycle;
\filldraw[fill=orange!20,draw=black]
  (3,0) rectangle (5,1);
\draw[dashed] (6,0) circle (5mm);
\end{lstlisting}
\texttt{cycle} menutup path ke titik awal. Urutan menggambar menentukan elemen mana yang tampak di depan.
\end{frame}
\begin{frame}[fragile]{Teks dan notasi matematika}
\begin{lstlisting}
\node at (0,0) {Titik asal};
\node[above] at (2,1) {$P=(2,1)$};
\node[draw,rounded corners,align=center] at (4,0)
  {Dua baris\\teks};
\end{lstlisting}
Gunakan \texttt{above}, \texttt{below}, \texttt{left}, atau \texttt{right} untuk menempatkan label relatif terhadap koordinat.
\end{frame}
\begin{frame}{Cek pemahaman: prediksi sebelum kompilasi}
\begin{enumerate}
\item Apa perbedaan garis \texttt{(0,0) -- +(1,0) -- +(0,1)} dan versi dengan \texttt{++}?
\item Apa yang terjadi jika titik koma setelah \texttt{\textbackslash draw} hilang?
\item Apa perbedaan \texttt{\textbackslash draw}, \texttt{\textbackslash fill}, dan \texttt{\textbackslash filldraw}?
\end{enumerate}
Diskusikan dengan pasangan selama 3 menit; uji prediksi di editor.
\end{frame}
\begin{frame}{Latihan 1A: ilustrasi dasar (10 menit)}
\begin{enumerate}
\item Gambar sumbu $x$ dan $y$ dengan panah serta label.
\item Tambahkan sebuah lingkaran dan persegi panjang berwarna.
\item Beri label teks dan satu notasi matematika.
\item Tambahkan garis putus-putus atau koordinat relatif.
\end{enumerate}
Kriteria: dokumen dikompilasi, bentuk dan label terlihat, serta perintah dapat dijelaskan.
\end{frame}
\begin{frame}{Diagnosis: periksa kesalahan}
Kerjakan lima potongan salah di \texttt{latihan.md}; baca kesalahan pertama pada log dan perbaiki satu hal setiap kali.
\medskip

Ubah warna atau bentuk dan jelaskan baris yang diubah. Bedakan \texttt{draw} dan \texttt{fill}. Bandingkan hasil dengan \texttt{contoh-sesi-1.tex}.
\end{frame}
\begin{frame}{Sesi 1: dari bentuk ke diagram terstruktur}
Gunakan node bernama dan gaya untuk membangun hubungan yang tetap mudah diubah. Selanjutnya kita pelajari positioning, scope, pengelompokan, dan perulangan.
\end{frame}
\begin{frame}[fragile]{Node bernama, anchor, dan panah}
\begin{lstlisting}
\node[draw] (a) at (0,0) {Masalah};
\node[draw] (b) at (3,0) {Metode};
\draw[->] (a.east) -- node[above] {dipelajari} (b.west);
\end{lstlisting}
Nama node memungkinkan hubungan tetap jelas saat teks berubah. Anchor seperti \texttt{north}, \texttt{south}, \texttt{east}, dan \texttt{west} menunjukkan titik pada batas node.
\end{frame}
\begin{frame}{Mengapa struktur kode penting?}
Bayangkan satu diagram dengan 12 node dan semua ukuran ditulis berulang. Perubahan desain mengharuskan banyak edit dan rawan tidak konsisten.
\medskip

Gunakan nama node untuk relasi, gaya untuk tampilan berulang, dan positioning untuk jarak antarnode. Uji dengan mengganti label menjadi lebih panjang.
\end{frame}
\begin{frame}[fragile]{Gaya dan positioning}
\begin{lstlisting}
\usetikzlibrary{positioning,arrows.meta}
\tikzset{tahap/.style={draw,rounded corners,
  fill=blue!10,align=center,minimum width=2cm}}
\node[tahap] (a) {Rencana};
\node[tahap,right=1cm of a] (b) {Analisis};
\draw[-{Stealth}] (a) -- (b);
\end{lstlisting}
Ubah definisi gaya sekali untuk memperbarui semua node yang memakainya.
\end{frame}
\begin{frame}[fragile]{Scope, perulangan, dan pengelompokan}
\begin{lstlisting}
\begin{scope}[xshift=2cm,rotate=10]
  \draw (0,0) -- (1,0);
\end{scope}
\foreach \x in {0,1,2} {\fill (\x,0) circle (2pt);}
% dengan pustaka fit:
\node[draw,fit=(a)(b),inner sep=4mm] {};
\end{lstlisting}
Latihan: alur bercabang dengan gaya, kelompok, dan elemen berulang. Periksa saat teks node berubah.
\end{frame}
\begin{frame}{Kualitas diagram}
Apakah arah relasi benar? Apakah teks terbaca pada ukuran akhir? Apakah jarak dan perataan konsisten? Gunakan bentuk, label, atau pola selain warna untuk membedakan kategori.
\end{frame}
\begin{frame}[fragile]{Posisi relatif dan jarak}
\begin{lstlisting}
\usetikzlibrary{positioning}
\begin{tikzpicture}[node distance=1cm and 2cm]
\node[draw] (a) {Mulai};
\node[draw,right=of a] (b) {Proses};
\node[draw,below=of b] (c) {Selesai};
\draw[->] (a) -- (b) -- (c);
\end{tikzpicture}
\end{lstlisting}
\texttt{node distance} mengatur jarak vertikal dan horizontal antartepi node. Periksa perbedaan saat label node diperpanjang.
\end{frame}
\begin{frame}[fragile]{Refaktor: pisahkan gaya dan isi}
\begin{lstlisting}
\tikzset{
  tahap/.style={draw,rounded corners,fill=blue!10,
    minimum width=2.5cm,minimum height=8mm,align=center},
  alur/.style={-{Stealth},thick}
}
\node[tahap] (a) {Pertanyaan};
\node[tahap,right=of a] (b) {Data};
\draw[alur] (a) -- (b);
\end{lstlisting}
Tugas singkat: ubah warna dan ketebalan melalui definisi gaya, lalu hitung berapa baris yang perlu diedit.
\end{frame}
\begin{frame}[fragile]{Percabangan dan jalur alternatif}
\begin{lstlisting}
\usetikzlibrary{shapes.geometric,positioning,arrows.meta}
\node[draw,diamond,aspect=2] (cek) {Cukup?};
\node[draw,above right=of cek] (ya) {Analisis};
\node[draw,below right=of cek] (tidak) {Tambah data};
\draw[-{Stealth}] (cek) -- node[above] {Ya} (ya);
\draw[-{Stealth}] (cek) -- node[below] {Tidak} (tidak);
\end{lstlisting}
Label cabang memberi makna yang tidak bergantung pada warna.
\end{frame}
\begin{frame}[fragile]{Scope dan transformasi}
\begin{lstlisting}
\draw[help lines] (0,0) grid (4,2);
\begin{scope}[xshift=2cm,rotate=30]
  \draw[blue,thick,->] (0,0) -- (1,0);
\end{scope}
\draw[red,->] (0,0) -- (1,0);
\end{lstlisting}
Transformasi dan opsi di dalam scope bersifat lokal. Prediksi arah dan lokasi kedua panah sebelum mengompilasi.
\end{frame}
\begin{frame}[fragile]{Perulangan dan kelompok}
\begin{lstlisting}
\foreach \x in {0,1,2,3}
  {\fill (0.4*\x,0) circle (2pt);}
% Pustaka: fit, backgrounds
\begin{scope}[on background layer]
  \node[draw,dashed,fit=(a)(b),inner sep=4mm] {};
\end{scope}
\end{lstlisting}
Loop cocok untuk elemen berulang; \texttt{fit} membungkus node yang sudah didefinisikan. Background menjaga kotak di belakang node.
\end{frame}
\begin{frame}{Latihan 1B: alur penelitian (15 menit)}
Buat minimal tiga tahap, satu cabang, dua gaya, dan node bernama. Jika waktu cukup, tambahkan kelompok dengan \texttt{fit} dan elemen berulang dengan \texttt{\textbackslash foreach}.
\medskip

Uji dengan mengganti salah satu node menjadi dua baris. Apakah panah dan jarak masih masuk akal? Gunakan \texttt{contoh-sesi-2.tex} untuk membandingkan pendekatan.
\end{frame}
\begin{frame}{Review dan exit ticket}
\begin{enumerate}
\item Temukan satu contoh kode berulang dan ubah menjadi gaya.
\item Periksa panah, label cabang, perataan, dan kontras.
\item Jelaskan apa yang berubah ketika sebuah gaya diubah.
\end{enumerate}
Tantangan: tambahkan panah balik tanpa menutupi teks.
\end{frame}
\section{Sesi 2: Grafik akademik, integrasi, dan proyek}
\begin{frame}{Sesi 2: tujuan dan alokasi 75 menit}
\begin{itemize}
\item 0--15: prinsip grafik dan demonstrasi kurva serta data.
\item 15--35: latihan grafik, legenda, dan anotasi.
\item 35--45: integrasi gambar pada dokumen dan Beamer.
\item 45--60: adaptasi gambar untuk bahan ajar atau penelitian.
\item 60--70: ulasan sejawat dan revisi; 70--75: penyerahan.
\end{itemize}
Target: grafik berlabel serta contoh gambar siap digunakan dalam dokumen atau presentasi.
\end{frame}
\begin{frame}{Diagram dan grafik menjawab pertanyaan berbeda}
TikZ membantu menyatakan struktur dan hubungan; PGFPlots membangun sumbu serta memetakan nilai numerik ke posisi.
\medskip

Sebelum memplot, tentukan variabel, satuan, sumber data, audiens, dan pesan. Data dalam paket latihan bersifat fiktif dan hanya digunakan untuk demonstrasi.
\end{frame}
\begin{frame}[fragile]{Kurva dengan PGFPlots}
\begin{lstlisting}
\usepackage{pgfplots}
\pgfplotsset{compat=1.18}
\begin{tikzpicture}
\begin{axis}[xlabel={$x$},ylabel={$f(x)$},grid=major]
  \addplot[blue,domain=0:5] {x^2};
\end{axis}
\end{tikzpicture}
\end{lstlisting}
Cantumkan variabel, satuan bila ada, dan domain yang sesuai.
\end{frame}
\begin{frame}[fragile]{Data, legenda, dan anotasi}
\begin{lstlisting}
\addplot[mark=*] table[x=x,y=y,col sep=comma] {data.csv};
\addlegendentry{Pengamatan}
\node[anchor=south east] at (axis cs:5,21) {Akhir};
\end{lstlisting}
Latihan: bandingkan data fiktif dengan kurva acuan. Periksa sumber, satuan, skala, dan posisi legenda. Lihat \texttt{contoh-sesi-3.tex}.
\end{frame}
\begin{frame}[fragile]{Kontrol domain, sampel, dan batas sumbu}
\begin{lstlisting}
\begin{axis}[
  xlabel={Waktu (minggu)}, ylabel={Nilai (satuan)},
  xmin=0,xmax=5,ymin=0,ymax=24,grid=major
]
\addplot[blue,domain=0:5,samples=50] {10+2*x};
\end{axis}
\end{lstlisting}
\texttt{domain} menentukan rentang fungsi yang dihitung; batas sumbu menentukan bagian yang ditampilkan. Bandingkan perubahan \texttt{ymin=0} dengan \texttt{ymin=10}.
\end{frame}
\begin{frame}[fragile]{Dari tabel CSV ke grafik}
\begin{lstlisting}
% data.csv:
% x,y
% 0,10
% 1,13
\addplot[blue,mark=*]
  table[x=x,y=y,col sep=comma] {data.csv};
\addlegendentry{Pengamatan fiktif}
\end{lstlisting}
Nama kolom harus sesuai header. Kompilasi dari folder yang memuat CSV atau sesuaikan jalurnya. Pastikan data dan sumber tersedia untuk reproduksi.
\end{frame}
\begin{frame}[fragile]{Bandingkan model dan pengamatan}
\begin{lstlisting}
\addplot[blue,mark=*] table[x=x,y=y,
  col sep=comma] {data.csv};
\addlegendentry{Pengamatan}
\addplot[black,dashed,domain=0:5,samples=2] {10+2*x};
\addlegendentry{Acuan}
\end{lstlisting}
Gunakan marker dan garis putus-putus untuk membedakan seri tanpa mengandalkan warna. Kurva acuan bukan bukti hubungan sebab akibat.
\end{frame}
\begin{frame}[fragile]{Anotasi dan koordinat sumbu}
\begin{lstlisting}
\node[anchor=south east,align=right]
  at (axis cs:5,21) {Titik terakhir};
\end{lstlisting}
\texttt{axis cs} menggunakan koordinat data dalam sistem sumbu. Tempatkan anotasi agar tidak menutupi titik, garis, atau legenda. Tambahkan caption yang menyatakan konteks dan sumber.
\end{frame}
\begin{frame}{Diskusi: apa yang dapat menyesatkan?}
\begin{itemize}
\item Sumbu tanpa satuan atau tanpa penjelasan variabel.
\item Pemotongan sumbu yang memperbesar kesan perbedaan.
\item Menghubungkan pengamatan yang tidak tepat diperlakukan sebagai kontinu.
\item Grafik tanpa sumber, periode, atau penjelasan data.
\end{itemize}
Skala tidak selalu harus mulai dari nol: pilih sesuai tujuan dan tunjukkan batas dengan jelas.
\end{frame}
\begin{frame}{Latihan 2: grafik dan diskusi (20 menit)}
Plot \texttt{data.csv} terhadap waktu; tambahkan acuan $10+2x$, legenda, anotasi titik terakhir, dan caption. Cek nilai dengan membaca tabel.
\medskip

Bandingkan dua batas sumbu $y$. Jelaskan apa yang berubah secara visual dan apa yang tetap sama pada datanya. Tantangan: masukkan data Anda sendiri dan cantumkan sumbernya.
\end{frame}
\begin{frame}[fragile]{Gambar dalam dokumen akademik}
\begin{lstlisting}
\begin{figure}[ht]
  \centering
  \begin{tikzpicture}
    % gambar Anda
  \end{tikzpicture}
  \caption{Hubungan konsep dalam penelitian.}
  \label{fig:konsep}
\end{figure}
Lihat Gambar~\ref{fig:konsep}.
\end{lstlisting}
Letakkan \texttt{\textbackslash label} setelah caption; kompilasi ulang untuk referensi silang.
\end{frame}
\begin{frame}{Merancang pengungkapan bertahap}
Dalam presentasi, ungkap elemen menurut urutan penjelasan: pertanyaan, metode, lalu relasi ke hasil.
\medskip

Pastikan setiap tahap masih bermakna, teks tidak berpindah secara mengganggu, dan slide akhir tetap terbaca. Uji seluruh overlay dalam PDF.
\end{frame}
\begin{frame}[fragile]{Presentasi bertahap di Beamer}
\begin{lstlisting}
\only<2->{\draw[->] (a) -- (b);}
\only<3->{\draw[->] (b) -- (c);}
\end{lstlisting}
Setiap overlay menambahkan relasi. Lihat \texttt{contoh-sesi-4.tex}, contoh Beamer untuk sesi 2.
\end{frame}
\begin{frame}[fragile]{Overlay: ruang dan urutan}
\begin{lstlisting}
% Node a, b, c sudah didefinisikan dalam tikzpicture
\only<2->{\draw[->] (a) -- (b);}
\only<3->{\draw[->] (b) -- (c);}
\end{lstlisting}
\texttt{\textbackslash only} tidak menempatkan isi pada overlay lain; perubahan batas gambar dapat menggeser letak. Uji PDF tiap tahap; pertahankan batas gambar atau gunakan mekanisme yang mempertahankan ruang bila diperlukan.
\end{frame}
\begin{frame}{Perencanaan proyek: lima menit}
\begin{enumerate}
\item Tulis satu kalimat pesan yang hendak disampaikan.
\item Pilih pembaca dan media: artikel, handout, atau proyektor.
\item Sketsa elemen dan relasinya; identifikasi apa yang perlu ditekankan.
\item Pilih titik awal: salah satu berkas contoh dalam paket.
\end{enumerate}
Periksa kebenaran isi sebelum memperhalus gaya.
\end{frame}
\begin{frame}{Proyek akhir}
Buat atau adaptasi gambar untuk satu materi kuliah atau penelitian. Sertakan caption atau legenda; gunakan gaya, node bernama, anotasi, dan scope jika sesuai dengan jenis gambar.
\medskip

Tukar dengan pasangan: periksa ketepatan isi, kompilasi, keterbacaan, dan kemudahan mengubah kode. Revisi lalu kumpulkan sumber dan PDF.
\end{frame}
\begin{frame}{Kriteria penilaian: 20 poin}
Lima kriteria, masing-masing 0--4: ketepatan isi; kompilasi dan kelengkapan; penggunaan konsep TikZ; keterbacaan; dan pemeliharaan kode.
\medskip

Nilai 4 berarti benar dan konsisten; 3 baik dengan kekurangan kecil; 2 sebagian benar; 1 perlu banyak perbaikan; 0 belum ada. Panduan lengkap dan target kelulusan ada di \texttt{panduan-pengajar.md}.
\end{frame}
\begin{frame}{Umpan balik sejawat: 10 menit}
\begin{enumerate}
\item Jelaskan pesan gambar dalam satu kalimat tanpa bantuan pembuatnya.
\item Periksa akurasi relasi atau data dan bagian yang sulit dibaca.
\item Ubah satu label atau gaya untuk menguji kemudahan pemeliharaan.
\item Beri satu kekuatan dan satu saran yang dapat ditindaklanjuti.
\end{enumerate}
Pembuat memilih revisi paling penting dan mengompilasi kembali.
\end{frame}
\begin{frame}{Presentasi dan penyerahan}
Tunjukkan gambar selama satu menit: jelaskan konteks, pilihan desain, dan satu perubahan berdasarkan umpan balik.
\medskip

Kumpulkan \texttt{.tex}, PDF, dan data atau berkas pendukung yang diperlukan. Jelaskan secara singkat apa yang ingin dikembangkan berikutnya.
\end{frame}
\begin{frame}{Pemeriksaan akhir}
\begin{enumerate}
\item Apakah gambar menyampaikan satu pesan yang jelas dan benar?
\item Apakah label, panah, skala, dan satuan dapat dipahami?
\item Apakah teks terbaca dan kategori tetap berbeda dalam cetak abu-abu?
\item Apakah kode dapat dikompilasi ulang bersama data yang diperlukan?
\end{enumerate}
\end{frame}
\begin{frame}{Pengayaan opsional: peta kuliah}
Contoh berikut dipertahankan sebagai bahan pengayaan atau pilihan proyek. Bagian ini berada di luar alokasi inti 75 menit.
\end{frame}
\begin{frame}{Dari tutorial ke peta yang dapat diajarkan}
Tutorial \href{https://tikz.dev/tutorial-map}{A Lecture Map for Johannes} menyusun konsep, anotasi kuliah, latar, dan kalender.
\medskip

Di sini kita mempertahankan empat fungsi itu dengan perintah yang telah dipelajari: node, koordinat, garis, gaya, \texttt{positioning}, \texttt{fit}, dan \texttt{scope}. Jadwal berupa daftar pertemuan, bukan kalender otomatis.
\end{frame}
\begin{frame}{Rancang struktur sebelum menulis kode}
\begin{enumerate}
\item Pilih satu mata kuliah dan tiga konsep utama.
\item Tentukan subkonsep dan relasi lintas cabang.
\item Pilih dua pertemuan, tanggal, dan hasil belajar yang terukur.
\item Tandai mana relasi hierarki dan mana relasi tambahan.
\end{enumerate}
Peta konsep bukan jadwal: satu pertemuan dapat membahas beberapa konsep dan satu konsep dapat muncul pada beberapa pertemuan.
\end{frame}
\begin{frame}[fragile]{Langkah 1: definisikan gaya}
\begin{lstlisting}
\usetikzlibrary{positioning,arrows.meta,fit,backgrounds}
\tikzset{
  konsep/.style={draw,rounded corners,fill=blue!10,
    minimum width=2.2cm,align=center},
  relasi/.style={-{Stealth},thick},
  catatan/.style={draw,fill=white,align=left}
}
\end{lstlisting}
Gaya memakai pola yang sudah dipraktikkan pada sesi 1.
\end{frame}
\begin{frame}[fragile]{Langkah 2: susun konsep dan relasi}
\begin{lstlisting}
\node[konsep] (pusat) {Metode penelitian};
\node[konsep,above right=of pusat] (tanya) {Pertanyaan};
\node[konsep,right=of pusat] (data) {Data};
\node[konsep,below right=of pusat] (analisis) {Analisis};
\draw[relasi] (pusat) -- (tanya);
\draw[relasi] (pusat) -- (data);
\draw[relasi] (pusat) -- (analisis);
\draw[relasi,dashed] (data) -- (analisis);
\end{lstlisting}
Garis putus-putus menandai relasi lintas cabang; jelaskan artinya pada legenda.
\end{frame}
\begin{frame}[fragile]{Langkah 3: anotasi pertemuan}
\begin{lstlisting}
\node[catatan,right=of tanya] (kuliah1)
  {P1: Merumuskan pertanyaan\\
   Target: menyusun pertanyaan riset};
\draw[relasi] (kuliah1.west) -- (tanya.east);
\end{lstlisting}
Gunakan target yang dapat diamati. Dalam contoh lengkap, jarak anotasi diatur agar tidak menutupi relasi lain.
\end{frame}
\begin{frame}[fragile]{Langkah 4: latar dan jadwal}
\begin{lstlisting}
\begin{scope}[on background layer]
  \node[draw,fill=blue!5,fit=(pusat)(tanya)(data)(analisis),
    inner sep=5mm] {};
\end{scope}
\node[catatan,below=of pusat,align=left]
  {Jadwal (contoh)\\P1: 3 Mar\\P2: 10 Mar};
\end{lstlisting}
Latar diletakkan di belakang; tanggal adalah contoh dan dapat diubah. Peta lengkap ada di \texttt{peta-kuliah.tex}.
\end{frame}
\begin{frame}{Pengayaan: peta kuliah sendiri}
Buat satu peta dengan konsep pusat, tiga konsep, dua subkonsep, satu hubungan lintas cabang, dua anotasi target belajar, latar kelompok, serta jadwal dua pertemuan.
\medskip

Gunakan paling sedikit dua gaya dan node bernama. Periksa keterbacaan, ruang antarelement, serta hubungan antara anotasi, konsep, dan jadwal.
\end{frame}
\begin{frame}{Ulasan dan exit ticket}
Tukar peta: apakah hubungan konsep benar, target dapat diamati, dan setiap jadwal konsisten dengan anotasinya? Berikan satu saran dan revisi.
\medskip

Exit ticket: jelaskan satu keputusan desain dan tunjukkan perintah yang menghubungkan dua konsep. Kumpulkan sumber dan PDF.
\end{frame}
% Sesi 3: materi belum ditentukan.
% Sesi 4: peta kuliah tersedia di folder ../Sesi 4.
\end{document}
